% Доклад по статье "Counting the Active Dimension of a Training Run:
% What a Single Scalar Log Can and Cannot Measure" (ICOMP 2026).
%
% Сборка: pdflatex main.tex (дважды). Кириллица через T2A + babel.
%
% Рисунки в figures/ порождаются скриптами fig_method.py, fig_validity.py,
% fig_setup.py, fig_grok.py, fig_backup.py из данных code/data/.
%
% Правила вёрстки, на которых всё держится.
%
% 1. Рисунок вставляется через \slidefig, то есть width=\textwidth и никакого
%    масштабирования. Каждый нарисован ровно на 5.945 in (= \textwidth), поэтому
%    кегль его подписей равен заданному в slide_style.py.
% 2. Место под текст берётся не из рисунка, а из объёма текста. Тело кадра
%    ~200 pt, рисунок 2.24 in = 161 pt, остаётся ~40 pt. При \small это ТРИ
%    строки, то есть около 230 знаков на весь слайд. Отсюда режим подписей:
%    только "куда смотреть" и одно число; разбор произносится вслух.
%    Проверять после каждой правки: `grep -c "Overfull .vbox" main.log`.
% 3. Заголовок кадра несёт вывод; тему повторять незачем.
% 4. Колонтитул показывает текущий раздел: в колоде 30 кадров, и без него
%    слушатель теряет, в какой части аргумента он находится.
\documentclass[aspectratio=169, 11pt]{beamer}

\usepackage{cmap}
\usepackage[T2A]{fontenc}
\usepackage[utf8]{inputenc}
\usepackage[english,russian]{babel}
\usepackage{amssymb,amsfonts,amsmath}
\usepackage{graphicx}

\usecolortheme{beaver}
\usefonttheme[onlymath]{serif}
\setbeamertemplate{navigation symbols}{}
\setbeamertemplate{blocks}[rounded=true,shadow=true]
\setbeamersize{text margin left=0.45cm,text margin right=0.45cm}
\setbeamerfont{author}{size=\normalsize}
\setbeamerfont{institute}{size=\normalsize}
\setbeamerfont{date}{size=\normalsize}
\setbeamerfont{frametitle}{size=\normalsize,series=\bfseries}
\setbeamertemplate{itemize items}[circle]

% Цвета совпадают с палитрой рисунков (Paul Tol high-contrast): один и тот же
% оттенок значит один и тот же режим на слайде и в статье.
\definecolor{recurrent}{HTML}{004488}
\definecolor{stochastic}{HTML}{BB5566}
\definecolor{transient}{HTML}{997700}
\definecolor{good}{HTML}{117733}
\setbeamercolor{structure}{fg=recurrent}
\setbeamercolor{alerted text}{fg=stochastic}
\setbeamercolor{block title}{bg=recurrent!14, fg=recurrent}
\setbeamercolor{block body}{bg=recurrent!5, fg=black}
\setbeamercolor{title}{fg=recurrent}
\setbeamercolor{frametitle}{fg=recurrent, bg=recurrent!10}

% Колонтитул: раздел слева, номер справа. Один тонкий ряд, 6 pt высоты.
\setbeamercolor{footline}{fg=recurrent!70!white}
\setbeamertemplate{footline}{%
  \hbox{\begin{beamercolorbox}[wd=\paperwidth,ht=1.7ex,dp=0.7ex,
        leftskip=0.45cm,rightskip=0.45cm]{footline}%
    \usebeamerfont{footline}\scriptsize\insertsectionhead\hfill
    \insertframenumber\,/\,\inserttotalframenumber%
  \end{beamercolorbox}}%
}

% Математические обозначения статьи, чтобы формулы на слайдах и в статье
% совпадали буква в букву.
\newcommand{\gW}{\mathcal{W}}
\newcommand{\gM}{\mathcal{M}}
\newcommand{\gA}{\mathcal{A}}
\newcommand{\PR}{\operatorname{PR}}
\newcommand{\R}{\mathbb{R}}
\newcommand{\vc}{\mathbf{c}}
\newcommand{\vy}{\mathbf{y}}
\newcommand{\vs}{\mathbf{s}}
\newcommand{\vu}{\mathbf{u}}
\newcommand{\vtheta}{\boldsymbol{\theta}}
\newcommand{\mV}{\mathbf{V}}
\newcommand{\dact}{d_{\mathrm{act}}}
\newcommand{\dmg}{\hat{d}_{\mathrm{MG}}}
\newcommand{\tgen}{t_{\mathrm{gen}}}
\newcommand{\rident}{\rho_{\mathrm{ident}}}
\newcommand{\WT}{W_T}

\newcommand{\slidefig}[1]{\centerline{\includegraphics[width=\textwidth]{#1}}}

% "Куда смотреть": две-три строки под рисунком, кеглем \small. Больше не влезет.
\newenvironment{look}
  {\vspace{0.2mm}\small\begin{itemize}%
   \setlength\itemsep{0.4mm}\setlength\parskip{0pt}}
  {\end{itemize}}

\title[Активная размерность обучения]{Подсчёт активной размерности процесса
обучения:\\ что может и чего не может измерить один скалярный лог}
\author{
    Николай Карлов\\
    Консультант: [имя консультанта]\\
    Эксперт: [имя эксперта]
}
\institute{
    Кафедра интеллектуальных систем, МФТИ\\
    По материалам статьи на ICOMP 2026
}
\date{2026}

\begin{document}

\begin{frame}
    \thispagestyle{empty}
    \maketitle
\end{frame}

% =============================================================== постановка ===
\section{Постановка}

\begin{frame}{Активная размерность: неформальная постановка}
    \begin{block}{Что определяется}
        \small
        Возьмём участок обучения, на котором закон обновления не меняется, а
        траектория параметров возвращается в одну и ту же область. Рассмотрим
        множество $\gA$ тех точек, к которым траектория подходит сколь угодно
        близко бесконечно много раз.
        \\[0.7mm]
        \alert{\textbf{Активная размерность --- размерность множества $\gA$:
        число независимых координат, которыми точка этого множества задаётся
        однозначно.}} Эти координаты и называются \textbf{активными
        компонентами}. Формальные определения --- на следующих двух слайдах.
    \end{block}

    \begin{block}{Три величины, которые различаются}
        \small
        \begin{tabular}{@{}l@{\ ---\ }l@{}}
          $P$        & число параметров: у трансформера здесь $226\,816$;\\
          $\PR$      & $\bigl(\sum_i \lambda_i\bigr)^2 \big/ \sum_i \lambda_i^2$
                       по спектру $\{\lambda_i\}$ ковариации;\\
          $\dact$    & размерность множества $\gA$.\\
        \end{tabular}
        \\[0.5mm]
        Обе записи слайда~12 имеют ковариацию задержек ранга $8$. У первой
        $\dact = 1$ и $\PR = 4.1$, у второй $\dact = 4$ и $\PR = 2.2$:
        \alert{порядок $\PR$ обратен порядку $\dact$} на всех восьми сидах.
    \end{block}
\end{frame}

\begin{frame}{Определение строго: режим и его мера}
    \begin{block}{Состояние и координаты}
        \small
        Состояние --- пара $\vs_t = (\vtheta_t, \vu_t)$: параметры
        $\vtheta_t \in \R^P$ и внутренние переменные оптимизатора $\vu_t$ (у
        AdamW --- два накопленных момента, у градиентного спуска их нет).
        \\[0.7mm]
        Если параметры лежат в аффинном подпространстве
        $\vtheta = \vtheta_0 + \mV^{\!\top}\vc$, $\mV\mV^{\!\top} = I_k$, то
        переход к координатам $\vc \in \R^k$ --- изометрия и размерность не
        меняет; дальше пишем $\vc$, для трансформера $\vc = \vtheta$.
    \end{block}

    \begin{block}{Режим и мера занятости}
        \small
        \textbf{Режим} $R$ --- участок работы, на котором (i) закон обновления и
        внешние условия неизменны; (ii) у состояния существует инвариантная
        вероятностная мера $\nu_R$; (iii) запись успевает её семплировать. Для
        эргодической детерминированной системы $\nu_R(C)$ --- \textbf{доля
        времени в} $C$:
        \[
          \nu_R(C) \;=\; \lim_{T \to \infty} \frac1T
          \sum_{t < T} \mathbf 1\{\vs_t \in C\},
          \qquad \mathbf 1\{\cdot\} \ \text{--- индикатор.}
        \]
        \textbf{Мера занятости} --- образ $\nu_R$ при $\pi_c : \vs \mapsto \vc$:
        $\mu_R(B) := \nu_R\bigl(\pi_c^{-1}(B)\bigr)$, $B \subseteq \R^k$
        борелевское.
    \end{block}
\end{frame}

\begin{frame}{Определение строго: размерность и компоненты}
    % Само определение стоит без бокса и по центру: это то единственное, что
    % слушатель должен унести с этого кадра, а блок вокруг него делает его одним
    % из трёх равноправных.
    \vspace{-1.0mm}
    \[
      d_{\mu_R}(\vc) \;:=\; \lim_{r \downarrow 0}
      \frac{\log \mu_R\bigl(B(\vc, r)\bigr)}{\log r},
      \qquad
      \alert{\dact(R) := d, \ \text{ если } \ d_{\mu_R}(\vc) = d \ \text{ для }
      \mu_R\text{-п.в. } \vc .}
    \]
    \vspace{-3.2mm}

    {\small Здесь $B(\vc, r) := \{\vc' : \|\vc' - \vc\| < r\}$. Мера с таким $d$
    называется \emph{точно размерной}; MG оценивает ровно этот показатель.}

    \smallskip

    \begin{block}{Гладкий случай: откуда берутся компоненты}
        \small
        Пусть $\gA_R := \operatorname{supp}\mu_R = \{\vc : \mu_R(B(\vc,r)) > 0\
        \ \forall r > 0\}$: точки, любой шар вокруг которых посещается с
        положительной долей времени. Если $\gA_R$ --- гладкое $d$-мерное
        многообразие, а плотность $\mu_R$ на нём положительна, то показатель
        равен $d$, и у почти всех точек есть карта
        $\Psi_R : U \subseteq \R^{d} \to \gA_R$ с
        $\operatorname{rank} D\Psi_R = d$. Её координаты $z_1, \dots, z_d$ ---
        \textbf{активные компоненты}; другая карта даёт другие,
        \alert{их число $d$ то же}.
    \end{block}

    \begin{block}{Роль окна}
        \small
        Окно $\gW$ из $L$ шагов задаёт
        $\widehat\mu_{\gW} := \tfrac1L \sum_{t \in \gW} \delta_{\vc_t}$, где
        $\delta_{\vc}$ --- единичная масса в $\vc$. Её носитель --- $L$ точек,
        показатель равен нулю: окно даёт \textbf{выборку} из $\mu_R$, по которой
        оценивают $\dact(R)$.
    \end{block}
\end{frame}

% =================================================================== метод ===
\section{Метод}

\begin{frame}{Шаг 1: из одного скаляра --- вектор задержек}
    \begin{block}{Наблюдение}
        \small
        Видим скаляр $x_t = \phi(\vs_t)$, где $\vs_t$ --- полное состояние
        оптимизатора, а $\phi$ фиксирована \emph{до} запуска. Строим
        \[
          \vy_t \;=\; \bigl(x_t,\ x_{t-\tau},\ \dots,\ x_{t-(E-1)\tau}\bigr)
          \;\in\; \R^{E}.
        \]
    \end{block}

    \begin{block}{Почему это законно}
        \small
        \textbf{Такенс, Сауэр и др.:} для \emph{типичной} пары (динамика,
        наблюдение) на компактном инвариантном множестве размерности $d$ это
        отображение --- вложение при $E > 2d$, то есть замена координат, не
        меняющая размерности.\\[0.6mm]
        Типичность для заранее выбранной $\phi$ \alert{не проверяема} --- отсюда
        все результаты сразу по семейству из 10 наблюдателей.
    \end{block}
\end{frame}

\begin{frame}{Шаг 2: показатель массы на конечном масштабе}
    \begin{block}{Конечные данные разрешают только диапазон радиусов}
        \small
        Предел $r \downarrow 0$ по конечной записи недостижим; для $0<r_1<r_2$
        \[
          d_R(\vc; r_1, r_2) \;:=\;
          \frac{\log\mu_R\bigl(B(\vc,r_2)\bigr)
                - \log\mu_R\bigl(B(\vc,r_1)\bigr)}
               {\log r_2 - \log r_1}.
        \]
        Компонента с амплитудой ниже $r_1$ \alert{не разрешается}.
    \end{block}

    \begin{block}{Оценка (MacKay--Ghahramani)}
        \small
        $r^{(i)}_1 \le \dots \le r^{(i)}_m$ --- расстояния от точки $i$ до её $m$
        соседей; сосед учитывается лишь при $|t_i - t_j| > \WT$ (исключение
        Тейлера). Пулинг правдоподобия по $N$ точкам \emph{до} обращения:
        \[
          \alert{\dmg^{\,-1}(\gW)} \;=\; \frac{1}{N(m-1)}
          \sum_{i=1}^{N} \sum_{j=1}^{m-1} \log \frac{r^{(i)}_m}{r^{(i)}_j}.
        \]
        MG оценивает \textbf{тот же показатель, но на масштабе самих расстояний
        до соседей}.
    \end{block}
\end{frame}

\begin{frame}{Когда оценку можно назвать размерностью}
    \begin{block}{Четыре условия}
        \small
        \begin{enumerate}\setlength\itemsep{0.3mm}
          \item окно целиком лежит в одном режиме $R$;
          \item динамика покрывает $\gA_R$ достаточно плотно --- \textbf{нужны
                возвраты};
          \item радиусы соседей лежат в устойчивом диапазоне скейлинга;
          \item delay map --- корректное вложение.
        \end{enumerate}
        Иначе сообщается \textbf{только статистика} $\dmg(\gW)$, без знака
        равенства. Заморожено: $\tau = 4$, $E \le 20$, $m = 20$, $\WT = 76$,
        $L = 8000$.
    \end{block}

    \begin{block}{Четыре величины, которые нельзя путать}
        \small
        \begin{tabular}{@{}ll@{}}
          $\dact(R)$ & размерность режима\\
          $\dact(R; r_1, r_2)$ & её конечномасштабная версия\\
          $\dmg(\gW)$ & вычисленная оконная статистика\\
          $\PR(\gW)$ & effective rank: линейный, вспомогательный\\
        \end{tabular}
    \end{block}
\end{frame}

\begin{frame}{Метод: скаляр $\to$ облако задержек $\to$ счёт соседей}
    \slidefig{figures/p_pipeline}
    \begin{look}
        \item Один и тот же запуск: $\dact = 1$, одна фаза. В \textbf{(b)} лежит
              \emph{реконструкция, сохранённая самим оценщиком}; здесь ничего не
              пересчитывается. Занятое множество --- замкнутая петля.
        \item Золотые в \textbf{(b)} --- сэмплы ближе $\WT$ по времени: период
              петли $\approx 25$ шагов, поэтому они разбросаны по всей петле, а
              соседи в \textbf{(c)} --- возвраты с других витков.
    \end{look}
\end{frame}

\begin{frame}{Выбор конкретного оценщика почти ничего не меняет}
    \slidefig{figures/p_estimator}
    \begin{look}
        \item Стенд: голова в 10-мерном подпространстве над замороженной сетью;
              данные модулированы синусоидами с несоизмеримыми частотами ---
              их число и есть истина.
        \item \textbf{(b)} $\rho$ --- ранговая корреляция с истиной. Держит
              \textbf{схема «задержки $+$ соседи»}; оценщик решает мало: линейный
              $\PR$ проигрывает лишь вдвое.
    \end{look}
\end{frame}

% ====================================================== что метод может ===
\section{Где оценка работает}

\begin{frame}{На построенной истине оценка попадает --- до потолка $\dact \approx 8$}
    \slidefig{figures/p_recovery}
    \begin{look}
        \item Ранги $1$--$8$; медиана по $10$ наблюдателям и $4$ сидам, полоса
              --- межквартильный разброс. (b) --- то же минус истины.
        \item Быстрый привод: промах $<0.7$ компоненты до $\dact = 8$, выше ---
              потолок конфигурации. В $25$ раз медленнее оценка встаёт на
              $3.4$. Ромб --- транзиент: $1$ против $16$.
    \end{look}
\end{frame}

\begin{frame}{Ошибка --- свойство наблюдателя не меньше, чем оценщика}
    \slidefig{figures/p_observers}
    \begin{look}
        \item Все $10$ наблюдателей, каждый --- фиксированная до запуска функция
              состояния оптимизатора. Ошибка на отложенных рангах $1,3,5,8$.
        \item Разброс впятеро: $0.34$ у фиксированной проекции параметров против
              $1.87$ у полнобатчевого лосса. Медиана $0.48$.
    \end{look}
\end{frame}

\begin{frame}{К самой геометрии оценка чувствительна, а шероховатость --- нет}
    \slidefig{figures/p_switch}
    \begin{look}
        \item \alert{Здесь нет сети:} два сигнала, каждый --- сумма четырёх
              синусоид. Частоты рационально независимы ($\dact = 4$) либо
              гармоники одной базы ($\dact = 1$, замыкание --- окружность).
        \item База второго \textbf{решена бисекцией} так, чтобы шероховатость
              совпала до последнего бита. Переключаем $4 \to 1 \to 4$; серые
              полосы --- смесь двух сигналов, истины там нет.
    \end{look}
\end{frame}

% ================================================== чего метод не может ===
\section{Где оценка не работает}

\begin{frame}{На входе: возвратный запуск и затухающий}
    \slidefig{figures/p_setup}
    \begin{look}
        \item \textbf{(a)} дрейф выключен: голову один раз толкнули из решения,
              $\vc^{*} + 0.8\,\mathbf{V}_{:r}\xi$, и она съезжает обратно.
              \textbf{100 \%} шагов вниз, ни одной смены знака.
        \item \textbf{(b)} дрейф включён: $1697$ смен знака. Невозвратный
              транзиент \alert{не семплирует инвариантную меру} --- $\gA_R$ у него
              нет. Сам путь --- кривая размерности $1$, но это другая величина.
    \end{look}
\end{frame}

\begin{frame}{«16» на затухающем запуске --- артефакт исключения}
    \slidefig{figures/p_theiler}
    \begin{look}
        \item Меняем \emph{только} $\WT$, всё остальное заморожено: на транзиенте
              статистика идёт от $1.20$ до $174$, на торе стоит. При $\WT = 0$
              она даёт размерность \emph{самого пути}; при $\WT = 76$ --- ничего.
        \item \textbf{(b)} почему: ближайшая точка у тора --- в $2554$ шагах, у
              транзиента --- в $6$. Запрет отнимает у одного всё, у другого
              ничего.
    \end{look}
\end{frame}

\begin{frame}{Шум входит в оценку: $2.5\,\%$ шума дают $11$ вместо $1$}
    \slidefig{figures/p_noise}
    \begin{look}
        \item Шум в $2.5\,\%$ от амплитуды дрейфа: при $\dact = 1$ оценка
              уезжает на $11.3$ и дальше почти не зависит от $r$.
        \item Пунктир --- истина \emph{чистого} дрейфа. У шумных ветвей
              occupation-мера может существовать, но её размерность считает
              \alert{и направления, возбуждённые шумом} --- это не та величина.
    \end{look}
\end{frame}

\begin{frame}{Режим виден по отношению $\rident$}
    \slidefig{figures/p_ident}
    \begin{look}
        \item Диагностика: считаем оценку при $E$ и при $2E$. Если это
              размерность множества, удвоение вложения её не меняет, и
              $\rident \approx 1$; если оценка следует за $E$, она уезжает.
        \item Два разных механизма шума --- гауссово форсирование шага и честный
              мини-батч --- дают \emph{один} вердикт $\rident \approx 1.55$.
              Транзиент уезжает ещё дальше.
    \end{look}
\end{frame}

% ================================================================ grokking ===
\section{Grokking}

\begin{frame}{У grokking'а оценка на дешёвом логе падает вдвое на переходе}
    \slidefig{figures/p_signal}
    \begin{look}
        \item Трансформер: 4 запуска с weight decay обобщают, 2 без него --- нет;
              всё выровнено по шагу генерализации $\tgen$.
        \item \textbf{(c)} по столбцу на запуск, сверху 4 обобщающих (розовые),
              снизу 2 без weight decay (золотые): глубина $\times 2$--$3$ при
              $p \le 0.075$ против $\times 1$ при $p \ge 0.95$.
    \end{look}
\end{frame}

\begin{frame}{Почему уровень читать нельзя 1: обе постановки вне области}
    \slidefig{figures/p_conf_map}
    \begin{look}
        \item Точка --- один лог. По горизонтали: сколько раз он пересёк линию
              тренда. По вертикали: $\rident$ (велико $=$ оценка следует за $E$).
        \item Все $10$ полнобатчевых логов пересекают тренд \textbf{ровно
              дважды} (полоса при $x = 2$). В синей области нет никого.
    \end{look}
\end{frame}

\begin{frame}{Почему уровень читать нельзя 2: падение бывает без обобщения}
    \slidefig{figures/p_conf_fall}
    \begin{look}
        \item \textbf{(a)} два запуска без weight decay: оценка на логе лосса
              падает впятеро.
        \item \textbf{(b)} а валидация у них так и остаётся на уровне шума,
              $0.08$--$0.09$ против $1.00$ у обобщающих.
    \end{look}
\end{frame}

\begin{frame}{Почему уровень читать нельзя 3: есть механизм без размерности}
    \slidefig{figures/p_conf_gain}
    \begin{look}
        \item Истина всё время одна, $r = 4$; меняем по одному фактору, который
              размерность не задевает. Серые почти ничего не двигают.
        \item Два розовых --- рост амплитуды и рост усиления лога --- дают
              $+1.5$ компоненты. Ровно это и делает растущая норма параметров.
    \end{look}
\end{frame}

\begin{frame}{Но это не ложная тревога: система там правда упрощается}
    \slidefig{figures/p_simplify}
    \begin{look}
        \item \textbf{(a)} у обоих запусков без weight decay норма растёт на
              \textbf{100 \%} логируемых шагов --- это и есть скатывание в тупое
              наращивание нормы.
        \item \textbf{(b)} у одного из них ранг пути ровно $1.00$: всё движение
              сводится к одному направлению. \alert{Оценка говорит «стало проще»
              --- и не врёт.}
    \end{look}
\end{frame}

% ================================================== прямое измерение ===
\section{Прямое измерение}

\begin{frame}{Независимый измеритель: CountSketch и эффективный ранг}
    \slidefig{figures/p_sketch}
    \begin{look}
        \item Каждый логируемый шаг сжимаем случайной проекцией в
              $\mathbb{R}^{1024}$; статистика --- эффективный ранг ковариации.
        \item \textbf{(a)} врёт не больше $0.11$; \textbf{(b)} памяти в $36$ раз
              меньше при $+6\,\%$ времени.
    \end{look}
\end{frame}

\begin{frame}{Прямое измерение: у перехода траектория сжимается и разжимается}
    \slidefig{figures/p_dip}
    \begin{look}
        \item Ранг коллапсирует у перехода почти до единицы и \textbf{снова
              расширяется} выше плато; у запусков без weight decay этого нет.
        \item \textbf{(c)} снимает возражение «траектория встала»: путь достигает
              дна позже и не возвращается.
    \end{look}
\end{frame}

\begin{frame}{Глубина падения запуски не разделяет}
    \slidefig{figures/p_depth}
    \begin{look}
        \item Во сколько раз упал эффективный ранг у перехода. У обобщающих
              $\times 2.2$--$\times 5.2$, у контролей $\times 1.0$--$\times 1.2$.
        \item Но засечка --- самое глубокое падение \emph{где угодно} в том же
              контрольном запуске: $2.37$, то есть не хуже, чем $2.24$ и $2.39$ у
              обобщающих. \alert{Одной глубины не хватает.}
    \end{look}
\end{frame}

\begin{frame}{Разделяют время падения и возврат}
    \slidefig{figures/p_timing}
    \begin{look}
        \item $\tgen$ у четырёх запусков различаются вчетверо ($3.7$--$13.7$ тыс.
              шагов), а минимум у каждого лежит в пределах $1.6$ тыс. шагов от
              своего $\tgen$ --- и в параметрах, и в функциях.
        \item Пунктир --- где лежал бы минимум, если бы коллапс случался
              на фиксированном абсолютном шаге. Точки на нём не лежат.
    \end{look}
\end{frame}

% =================================================================== итоги ===
\section{Итоги}

\begin{frame}{Итоги: сигнал об упрощении есть, размерности в нём нет}
    \small
    \begin{enumerate}\setlength\itemsep{1.5mm}
        \item \textbf{Где $\dmg(\gW)$ можно назвать $\dact(R)$.} В
              детерминированном и возвратном режиме с согласованным лагом --- до
              потолка $\approx 8$ компонент. Держит схема «задержки $+$
              соседи»; конкретный оценщик почти не важен.
        \item \textbf{В grokking'е --- нельзя.} Диагностики отвергают обе
              постановки; окно у перехода не обязано семплировать одну
              стационарную меру. Поэтому падение --- \alert{сигнал,
              скоррелированный с переходом}; измерением $\dact(R)$ его называть
              нельзя.
        \item \alert{\textbf{Но как детектор упрощения метод работает.}} И у
              обобщающих запусков, и без weight decay он ловит реальное событие: там
              --- временное сжатие у перехода, здесь --- скатывание в рост нормы
              с траекторией ранга 1. Ложных тревог нет; есть неразличение
              \emph{вида} упрощения.
        \item \textbf{Вид различает прямое измерение.} Сжатие у grokking'а
              временное и \emph{обращается}, у запуска без weight decay --- постоянное.
              Разделяют время и возврат; глубина --- нет.
    \end{enumerate}
\end{frame}

\begin{frame}{Ограничения, честно}
    \small
    \begin{itemize}\setlength\itemsep{1.6mm}
        \item Нужен возвратный режим и лаг, согласованный с периодом, которого лог
              не выдаёт. Выше $\approx 8$ компонент оценка только порядковая, и
              линейный $\mathrm{PR}$ по задержкам почти не хуже.
        \item Условие про \textbf{устойчивый диапазон скейлинга} мы не проверяем:
              $r_1, r_2$ задаются самими расстояниями до соседей, отбора области
              нет. Для шумного режима это очевидное улучшение протокола, для
              транзиента бесполезное --- там все двадцать расстояний лежат в
              полосе шириной 13 \%.
        \item Коллапс --- четыре запуска против двух, разделённых weight decay,
              который заодно двигает норму. Метка ставится только \emph{после}
              перехода.
        \item В полнобатчевой постановке результат не воспроизводится.
    \end{itemize}
\end{frame}

% ================================================================= резерв ===
\appendix
\section{Резерв}

\begin{frame}{Резерв: выбор лага и оцениваемая величина}
    \slidefig{figures/p_backup_tau}
    \begin{look}
        \item \textbf{(a)} каждый ранг поделён на свою истину. Окно задержки
              должно покрывать заметную долю периода; лог обучения периода не даёт.
        \item \textbf{(b)} при амплитудах фаз $q^l$ размерность остаётся $r$, а
              эффективный ранг падает. Оценка следует за \emph{размерностью}.
    \end{look}
\end{frame}

\begin{frame}{Резерв: суррогатный тест для падения на логе}
    \slidefig{figures/p_backup_surrogate}
    \begin{look}
        \item Запуск сравнивается \emph{сам с собой}: суррогаты хранят спектр, но
              рвут тонкую структуру. $p$ --- доля суррогатов, упавших не мельче.
        \item У двух запусков без weight decay на этом окне почти нет динамического
              диапазона, так что их $\times 1$ --- не дименсиональный негатив.
    \end{look}
\end{frame}

\begin{frame}{Резерв: четыре формы облака и оценка на каждой}
    \slidefig{figures/p_shapes}
    \begin{look}
        \item Плоскость первых двух delay-координат --- то множество, по которому
              идёт поиск соседей. У двух возвратных случаев оно заполнено, и
              оценка совпадает с истиной.
        \item У транзиента $\approx 16$: величина на пределе исключения,
              размерностью она не является.
    \end{look}
\end{frame}

\end{document}
